\documentclass[11pt]{article}

\usepackage[a4paper,margin=1in]{geometry}
\usepackage[T1]{fontenc}
\usepackage[utf8]{inputenc}
\usepackage{tabularx}
\usepackage{booktabs}
\usepackage{enumitem}
\usepackage{hyperref}
\usepackage{amsmath}
\usepackage{graphicx}
\usepackage{xcolor}
\usepackage{pifont}
\newcommand{\cmark}{\ding{51}}
\newcommand{\xmark}{\ding{55}}

\setlength{\parindent}{0pt}
\setlength{\parskip}{6pt}

\title{\textbf{Adaptive Subtitle Placement System}\\[4pt]
\large Technical Report}
\author{Adaptive Captioning Project}
\date{\today}

\begin{document}
\maketitle

\section*{Abstract}
This report describes an intelligent subtitle placement system for educational videos. The system automatically repositions subtitles across eight screen regions to avoid covering critical content such as code, equations, diagrams, and handwritten material. It uses computer vision to analyse each frame and select the optimal subtitle position, achieving \textbf{6.6$\times$ real-time} processing speed without modifying the original video.

\section{Problem Statement}
Standard subtitle systems place all captions at a fixed bottom-centre position. In educational videos this frequently covers:
\begin{itemize}[leftmargin=1.5em, itemsep=2pt]
    \item Code being written or demonstrated (programming tutorials)
    \item Mathematical equations and diagrams (STEM lectures)
    \item Whiteboard or smartboard content (handwritten teaching)
    \item PowerPoint slide text and figures (presentations)
\end{itemize}
This occlusion disrupts learning, particularly for students who depend on captions for accessibility or language support.

\section{System Overview}
The system takes a video file and an SRT subtitle file as input and produces an \textbf{adaptive WebVTT subtitle file} as output --- the original video is never modified.

\textbf{Processing pipeline (per subtitle):}
\begin{enumerate}[leftmargin=1.5em, itemsep=2pt]
    \item \textbf{Frame extraction} --- Three representative frames are decoded using GPU hardware acceleration (NVDEC).
    \item \textbf{Content analysis} --- Each of the 8 screen regions is scored using:
    \begin{itemize}
        \item \textit{OCR overlap} --- detects on-screen text (EasyOCR, GPU-accelerated)
        \item \textit{Edge density} --- detects diagrams, graphs, and dense visuals (Canny edge detection)
        \item \textit{Active writing} --- detects instructor writing live on a board
    \end{itemize}
    \item \textbf{Region selection} --- The region with the lowest combined score (least content) is chosen.
    \item \textbf{Temporal smoothing} --- Hysteresis and dwell-time logic prevent jittery position changes.
    \item \textbf{Output} --- Per-cue VTT position metadata is written for each subtitle entry.
\end{enumerate}

\section{Scoring Formula}
For each of the 8 regions, a score is computed. \textbf{Lower score = better placement.}

\[
\text{Score} = \underbrace{\text{base\_score}}_{\text{ergonomic preference}}
+ \; ({\text{ocr\_overlap} \times 5.0})
+ \; ({\text{edge\_density} \times 4.0})
+ \; ({\text{active\_writing} \times 3.0})
\]

Bottom regions carry a base score of 0 (preferred by default). Upper regions carry a penalty of up to 1.5 to discourage unnecessary upward displacement.

\subsection*{OCR Overlap (weight: 5.0)}
The frame is cropped to the region's bounding box and passed to EasyOCR (GPU-accelerated). The overlap is computed as:
\[
\text{ocr\_overlap} = \frac{\sum \text{detected text box areas}}{\text{region area}}
\]
A high value means the region contains significant on-screen text and should be avoided. Cropping before OCR rather than filtering full-frame results improves both accuracy and speed for small regions.

\subsection*{Edge Density (weight: 4.0)}
Canny edge detection is applied to the grayscale frame (thresholds: 50, 150). The edge map is computed \emph{once per frame} and reused across all 8 regions. The density for each region is:
\[
\text{edge\_density} = \frac{\text{edge pixels in region}}{\text{total region pixels}}
\]
This captures diagrams, graphs, geometric shapes, and dense slide content. If \(\text{ocr\_overlap} > 0.15\), edge density is set to zero to avoid double-penalising text-bearing regions.

\subsection*{Active Writing (weight: 3.0)}
This component detects whether an instructor is actively writing on a board during the subtitle window. A reference frame $f_0$ is taken at subtitle start $+\,0.1\text{s}$; subsequent frames are sampled every 10th frame. For each sample $f_i$:
\[
\text{diff}_i = |f_i^{\text{gray}} - f_0^{\text{gray}}|, \qquad
\text{changed mask} = \text{diff}_i > 0.06
\]
Connected component analysis rejects large blobs (slide transitions, presenter silhouette) and keeps only small, localised ink strokes. A region is flagged as active if ink persists in $\geq \lfloor n/3 \rfloor$ sampled frames. Only the top 2 most active regions are penalised (scores 1.0 and 0.67 respectively).

\section{Screen Regions}
The screen is divided into 8 subtitle placement regions:

\begin{center}
\begin{tabularx}{0.85\textwidth}{>{\bfseries}l c c}
\toprule
Region & Vertical Position & Base Score \\
\midrule
bottom\_center & Bottom 15\% & 0 (preferred) \\
bottom\_left / bottom\_right & Bottom 15\% & 0 \\
middle\_left / middle\_right & Middle band & 0.5 \\
upper\_left / upper\_right & Upper band & 1.0 \\
top\_center & Top 15\% & 1.5 (last resort) \\
\bottomrule
\end{tabularx}
\end{center}

\section{Performance}
Tested on a \textbf{14-minute educational video} (121 subtitles) on an NVIDIA RTX 2050 (mid-range laptop GPU):

\begin{center}
\begin{tabularx}{0.8\textwidth}{>{\bfseries}l r r}
\toprule
Metric & Before & After \\
\midrule
Total processing time & 1,625 s (27 min) & 125 s (2 min) \\
Real-time speed & 0.51$\times$ & \textbf{6.6$\times$} \\
Daily throughput & 2.2 h of video & 52.8 h of video \\
\bottomrule
\end{tabularx}
\end{center}

Key optimisations:
\begin{itemize}[leftmargin=1.5em, itemsep=2pt]
    \item \textbf{GPU hardware decoding} via FFmpeg NVDEC (14$\times$ speedup)
    \item \textbf{Edge map caching} --- Canny computed once per frame, reused for all 8 regions (54$\times$ speedup)
    \item \textbf{Similarity pre-filtering} --- Consecutive frames that are visually identical skip all analysis and inherit the previous subtitle region
\end{itemize}

\subsection*{Similarity Gate: Mean Absolute Deviation (MAD)}
Before running any OCR or edge analysis, the system checks whether the current frame is sufficiently similar to the previous one using three sequential tests. All three must pass for a frame to be considered a duplicate:

\begin{enumerate}[leftmargin=1.5em, itemsep=2pt]
    \item \textbf{Histogram correlation} --- A 64-pixel grayscale histogram is compared between frames. Frames must correlate at $\geq 0.92$ to proceed.
    \item \textbf{Mean Absolute Deviation (MAD)} --- The average per-pixel brightness difference is computed:
    \[
    \text{MAD} = \frac{1}{W \times H} \sum_{x,y} \left| f_{\text{current}}(x,y) - f_{\text{prev}}(x,y) \right|
    \]
    If $\text{MAD} < 3.0$ (out of 255), frames are considered visually identical.
    \item \textbf{Blurred difference} --- A Gaussian-blurred full-frame absolute difference is computed. If fewer than 100 pixels change significantly, the frame is marked as static.
\end{enumerate}

If all three conditions are satisfied, the subtitle inherits the previous region with zero re-computation. On slide-based videos this skips up to 95\% of entries, yielding a significant additional speedup.

\section{Validation}
Validated on \textbf{15 educational videos} across three content types:

\begin{center}
\begin{tabularx}{0.7\textwidth}{>{\bfseries}l c}
\toprule
Category & Count \\
\midrule
Smartboard / handwritten & 9 \\
PowerPoint presentations & 4 \\
Live coding / terminal & 2 \\
\bottomrule
\end{tabularx}
\end{center}

Quality check: adaptive vs.\ baseline outputs showed \textbf{99\%+ text content match} with only 3 of 617 VTT lines differing, confirming output correctness.

\section{Technology Stack}

\begin{center}
\begin{tabularx}{0.85\textwidth}{>{\bfseries}l X}
\toprule
Component & Technology \\
\midrule
Language & Python 3.11 \\
Deep Learning & PyTorch 2.5.1 + CUDA 12.1 \\
OCR & EasyOCR (GPU-accelerated) \\
Computer Vision & OpenCV 4.13 \\
Video Decoding & FFmpeg with NVDEC \\
Backend API & FastAPI \\
Frontend & Vite + React \\
\bottomrule
\end{tabularx}
\end{center}

\section{Conclusion}
The adaptive subtitle placement system successfully solves the content-occlusion problem in educational video subtitling. By combining OCR, edge detection, and active writing analysis with temporal smoothing, it achieves intelligent and stable subtitle positioning without modifying the source video. The system is production-ready, operating at 6.6$\times$ real-time speed on mid-range hardware, and has been validated across 15 diverse educational videos.

\end{document}